% --- Archon physics formalization source begin ---
% archon:physics
% archon:covers PhyXMiniProblems/problem_phyx_mini_0406.lean
% archon:source-report reports/phyx_mini/problem_phyx_mini_0406.source.json

\chapter{Physics problem phyx\_mini\_0406}
\label{ch:PhyXMiniProblems_problem_phyx_mini_0406}

\paragraph{Problem source.}
\#\# Physical scenario

A beaker with a metal bottom is filled with 20 \$	ext\{g\}\$ of water at 20 \$\textasciicircum{}\{\textbackslash{}circ\}\textbackslash{}text\{C\}\$. It is brought into good thermal contact with a 4000 \$	ext\{cm\}\textasciicircum{}3\$ container holding 0.40 \$	ext\{mol\}\$ of a monatomic gas at 10 \$	ext\{atm\}\$ pressure. Both containers are well insulated from their surroundings.You can assume that the containers themselves are nearly massless and do not affect the outcome.

\#\# Question

What is the gas pressure after a long time has elapsed?

\#\# Answer choices

- A: A: 5.0 atm

- B: B: 10.0 atm

- C: C: 2.8 atm

- D: D: 1.0 atm

\#\# Figure caption (auxiliary; use the image as primary evidence)

The image is a diagram showcasing a cross-sectional view of a container. It contains several labeled components:

1. **Water**: A blue section at the top, labeled "Water."
2. **Gas**: A larger yellow section beneath the blue section, labeled "Gas."
3. **Thin Metal**: A thin black line separating the water and gas sections, labeled "Thin metal."
4. **Insulation**: Surrounding both the water and gas sections is a gray area with small dots, labeled "Insulation."

The diagram visually distinguishes each material using color and texture, and labels indicate their respective positions and materials.

\#\# Dataset metadata

Ideal Gases and Kinetic Theory; Physical Model Grounding Reasoning, Multi-Formula Reasoning

\paragraph{Recorded answer/context.}
C: C: 2.8 atm

\paragraph{Figure/image path.}
/root/proposal\_for\_physic/science-mango/phyx\_mini\_run/phyx\_data/test\_image/406.png

\paragraph{Formalization target.}
create a compiling Lean file with sorry bodies at `PhyXMiniProblems/problem\_phyx\_mini\_0406.lean`. The Lean declarations must preserve the physical quantities, dimensions or dimensional roles, figure labels, governing-law hypotheses, and final relation expressed by this problem.
Use Mathlib/Physlib names found through LeanExplore where available. If a physics API is missing, introduce faithful local abstractions rather than scalar placeholder aliases.

\begin{theorem}[Physics formalization target]
\label{thm:physics:phyx_mini_0406:target}
\lean{PhyXMiniProblems.ProblemPhyXMini0406.problem_phyx_mini_0406}
\uses{def:physics:phyx-mini-0406:phyxminiproblems-problemphyxmini0406-amountofsubstancescale, def:physics:phyx-mini-0406:phyxminiproblems-problemphyxmini0406-specificheatcapacityscale, def:physics:phyx-mini-0406:phyxminiproblems-problemphyxmini0406-pressureinstandardatmospheres, def:physics:phyx-mini-0406:phyxminiproblems-problemphyxmini0406-insulatedgaswatersetup, def:physics:phyx-mini-0406:phyxminiproblems-problemphyxmini0406-matchesinsulatedgaswaterscenario, def:physics:phyx-mini-0406:phyxminiproblems-problemphyxmini0406-matchessuppliedthermalcontactfigure, def:physics:phyx-mini-0406:phyxminiproblems-problemphyxmini0406-matchesproblemreadouts, def:physics:phyx-mini-0406:phyxminiproblems-problemphyxmini0406-usesreferencewaterthermaldata, def:physics:phyx-mini-0406:phyxminiproblems-problemphyxmini0406-hasphysicalthermodynamicparameters, def:physics:phyx-mini-0406:phyxminiproblems-problemphyxmini0406-satisfiesinsulatedgaswaterlaws, def:physics:phyx-mini-0406:phyxminiproblems-problemphyxmini0406-recordeddatasetanswer, def:physics:phyx-mini-0406:phyxminiproblems-problemphyxmini0406-roundstodisplayedpressure, def:physics:phyx-mini-0406:phyxminiproblems-problemphyxmini0406-isuniqueclosestdisplayedpressure}
The assigned autoformalize agent should translate this physics problem into Lean declarations in the covered file, with theorem and lemma proof bodies written as `by sorry`.
\end{theorem}
\begin{proof}
This is an autoformalization task, not a proof task. Produce statements that can later be proved by the physics prover without weakening the physical model.
\end{proof}

% --- Archon declaration topology begin ---
\section{Declaration topology}
\begin{definition}[volume Dimension]
  \label{def:physics:phyx-mini-0406:phyxminiproblems-problemphyxmini0406-volumedimension}
  \lean{PhyXMiniProblems.ProblemPhyXMini0406.volumeDimension}
  The physical dimension of volume, L³
\end{definition}
\begin{proof}
  This is the declared model definition of volume Dimension.
\end{proof}
\begin{definition}[Mass Quantity]
  \label{def:physics:phyx-mini-0406:phyxminiproblems-problemphyxmini0406-massquantity}
  \lean{PhyXMiniProblems.ProblemPhyXMini0406.MassQuantity}
  A nonnegative, unit-independent physical mass
\end{definition}
\begin{proof}
  This is the declared model definition of Mass Quantity.
\end{proof}
\begin{definition}[Volume Quantity]
  \label{def:physics:phyx-mini-0406:phyxminiproblems-problemphyxmini0406-volumequantity}
  \lean{PhyXMiniProblems.ProblemPhyXMini0406.VolumeQuantity}
  \uses{def:physics:phyx-mini-0406:phyxminiproblems-problemphyxmini0406-volumedimension}
  A nonnegative, unit-independent physical volume
\end{definition}
\begin{proof}
  This is the declared model definition of Volume Quantity.
\end{proof}
\begin{definition}[Amount Of Substance Scale]
  \label{def:physics:phyx-mini-0406:phyxminiproblems-problemphyxmini0406-amountofsubstancescale}
  \lean{PhyXMiniProblems.ProblemPhyXMini0406.AmountOfSubstanceScale}
  An abstract physical amount-of-substance scale with a calibrated mole readout. The carrier is not identified with a bare real number
\end{definition}
\begin{proof}
  This is the declared model definition of Amount Of Substance Scale.
\end{proof}
\begin{definition}[Specific Heat Capacity Scale]
  \label{def:physics:phyx-mini-0406:phyxminiproblems-problemphyxmini0406-specificheatcapacityscale}
  \lean{PhyXMiniProblems.ProblemPhyXMini0406.SpecificHeatCapacityScale}
  An abstract scale for mass-specific heat capacity, calibrated in joules per kilogram-kelvin
\end{definition}
\begin{proof}
  This is the declared model definition of Specific Heat Capacity Scale.
\end{proof}
\begin{definition}[mass Readout]
  \label{def:physics:phyx-mini-0406:phyxminiproblems-problemphyxmini0406-massreadout}
  \lean{PhyXMiniProblems.ProblemPhyXMini0406.massReadout}
  \uses{def:physics:phyx-mini-0406:phyxminiproblems-problemphyxmini0406-massquantity}
  Read a physical mass in a selected mass unit
\end{definition}
\begin{proof}
  This is the declared model definition of mass Readout.
\end{proof}
\begin{definition}[volume Readout]
  \label{def:physics:phyx-mini-0406:phyxminiproblems-problemphyxmini0406-volumereadout}
  \lean{PhyXMiniProblems.ProblemPhyXMini0406.volumeReadout}
  \uses{def:physics:phyx-mini-0406:phyxminiproblems-problemphyxmini0406-volumequantity}
  Read a physical volume in the cube of a selected length unit
\end{definition}
\begin{proof}
  This is the declared model definition of volume Readout.
\end{proof}
\begin{definition}[pressure In Pascals]
  \label{def:physics:phyx-mini-0406:phyxminiproblems-problemphyxmini0406-pressureinpascals}
  \lean{PhyXMiniProblems.ProblemPhyXMini0406.pressureInPascals}
  Read a dimensionful pressure in coherent SI pascals
\end{definition}
\begin{proof}
  This is the declared model definition of pressure In Pascals.
\end{proof}
\begin{definition}[energy In Joules]
  \label{def:physics:phyx-mini-0406:phyxminiproblems-problemphyxmini0406-energyinjoules}
  \lean{PhyXMiniProblems.ProblemPhyXMini0406.energyInJoules}
  Read a dimensionful energy in coherent SI joules
\end{definition}
\begin{proof}
  This is the declared model definition of energy In Joules.
\end{proof}
\begin{definition}[mass In Grams]
  \label{def:physics:phyx-mini-0406:phyxminiproblems-problemphyxmini0406-massingrams}
  \lean{PhyXMiniProblems.ProblemPhyXMini0406.massInGrams}
  \uses{def:physics:phyx-mini-0406:phyxminiproblems-problemphyxmini0406-massquantity, def:physics:phyx-mini-0406:phyxminiproblems-problemphyxmini0406-massreadout}
  Gram readout used for the water mass in the problem statement
\end{definition}
\begin{proof}
  This is the declared model definition of mass In Grams.
\end{proof}
\begin{definition}[mass In Kilograms]
  \label{def:physics:phyx-mini-0406:phyxminiproblems-problemphyxmini0406-massinkilograms}
  \lean{PhyXMiniProblems.ProblemPhyXMini0406.massInKilograms}
  \uses{def:physics:phyx-mini-0406:phyxminiproblems-problemphyxmini0406-massquantity, def:physics:phyx-mini-0406:phyxminiproblems-problemphyxmini0406-massreadout}
  Kilogram readout used by the calorimetry law
\end{definition}
\begin{proof}
  This is the declared model definition of mass In Kilograms.
\end{proof}
\begin{definition}[volume In Cubic Centimeters]
  \label{def:physics:phyx-mini-0406:phyxminiproblems-problemphyxmini0406-volumeincubiccentimeters}
  \lean{PhyXMiniProblems.ProblemPhyXMini0406.volumeInCubicCentimeters}
  \uses{def:physics:phyx-mini-0406:phyxminiproblems-problemphyxmini0406-volumequantity, def:physics:phyx-mini-0406:phyxminiproblems-problemphyxmini0406-volumereadout}
  Cubic-centimetre readout used for the gas-container volume
\end{definition}
\begin{proof}
  This is the declared model definition of volume In Cubic Centimeters.
\end{proof}
\begin{definition}[volume In Cubic Meters]
  \label{def:physics:phyx-mini-0406:phyxminiproblems-problemphyxmini0406-volumeincubicmeters}
  \lean{PhyXMiniProblems.ProblemPhyXMini0406.volumeInCubicMeters}
  \uses{def:physics:phyx-mini-0406:phyxminiproblems-problemphyxmini0406-volumequantity, def:physics:phyx-mini-0406:phyxminiproblems-problemphyxmini0406-volumereadout}
  Cubic-metre readout used by the SI ideal-gas law
\end{definition}
\begin{proof}
  This is the declared model definition of volume In Cubic Meters.
\end{proof}
\begin{definition}[temperature In Kelvin]
  \label{def:physics:phyx-mini-0406:phyxminiproblems-problemphyxmini0406-temperatureinkelvin}
  \lean{PhyXMiniProblems.ProblemPhyXMini0406.temperatureInKelvin}
  Kelvin readout of Physlib's absolute temperature
\end{definition}
\begin{proof}
  This is the declared model definition of temperature In Kelvin.
\end{proof}
\begin{definition}[temperature In Degrees Celsius]
  \label{def:physics:phyx-mini-0406:phyxminiproblems-problemphyxmini0406-temperatureindegreescelsius}
  \lean{PhyXMiniProblems.ProblemPhyXMini0406.temperatureInDegreesCelsius}
  \uses{def:physics:phyx-mini-0406:phyxminiproblems-problemphyxmini0406-temperatureinkelvin}
  Celsius readout obtained from the absolute kelvin readout
\end{definition}
\begin{proof}
  This is the declared model definition of temperature In Degrees Celsius.
\end{proof}
\begin{definition}[pressure In Standard Atmospheres]
  \label{def:physics:phyx-mini-0406:phyxminiproblems-problemphyxmini0406-pressureinstandardatmospheres}
  \lean{PhyXMiniProblems.ProblemPhyXMini0406.pressureInStandardAtmospheres}
  \uses{def:physics:phyx-mini-0406:phyxminiproblems-problemphyxmini0406-pressureinpascals}
  Pressure as a multiple of Physlib's standard atmosphere
\end{definition}
\begin{proof}
  This is the declared model definition of pressure In Standard Atmospheres.
\end{proof}
\begin{definition}[universal Molar Gas Constant J per mol K]
  \label{def:physics:phyx-mini-0406:phyxminiproblems-problemphyxmini0406-universalmolargasconstant-j-per-mol-k}
  \lean{PhyXMiniProblems.ProblemPhyXMini0406.universalMolarGasConstant_J_per_mol_K}
  Universal molar gas constant, read in joules per mole-kelvin
\end{definition}
\begin{proof}
  This is the declared model definition of universal Molar Gas Constant J per mol K.
\end{proof}
\begin{definition}[Thermal Stage]
  \label{def:physics:phyx-mini-0406:phyxminiproblems-problemphyxmini0406-thermalstage}
  \lean{PhyXMiniProblems.ProblemPhyXMini0406.ThermalStage}
  The given initial state and the equilibrium state reached after a long time
\end{definition}
\begin{proof}
  This is the declared model definition of Thermal Stage.
\end{proof}
\begin{definition}[Gas Constitution]
  \label{def:physics:phyx-mini-0406:phyxminiproblems-problemphyxmini0406-gasconstitution}
  \lean{PhyXMiniProblems.ProblemPhyXMini0406.GasConstitution}
  The constitution of the gas named in the prose
\end{definition}
\begin{proof}
  This is the declared model definition of Gas Constitution.
\end{proof}
\begin{definition}[Interface Material]
  \label{def:physics:phyx-mini-0406:phyxminiproblems-problemphyxmini0406-interfacematerial}
  \lean{PhyXMiniProblems.ProblemPhyXMini0406.InterfaceMaterial}
  Material forming the thermally conducting interface
\end{definition}
\begin{proof}
  This is the declared model definition of Interface Material.
\end{proof}
\begin{definition}[Thermal Contact Quality]
  \label{def:physics:phyx-mini-0406:phyxminiproblems-problemphyxmini0406-thermalcontactquality}
  \lean{PhyXMiniProblems.ProblemPhyXMini0406.ThermalContactQuality}
  Quality of thermal contact across the beaker bottom
\end{definition}
\begin{proof}
  This is the declared model definition of Thermal Contact Quality.
\end{proof}
\begin{definition}[Exterior Boundary Condition]
  \label{def:physics:phyx-mini-0406:phyxminiproblems-problemphyxmini0406-exteriorboundarycondition}
  \lean{PhyXMiniProblems.ProblemPhyXMini0406.ExteriorBoundaryCondition}
  Thermal condition of the apparatus boundary
\end{definition}
\begin{proof}
  This is the declared model definition of Exterior Boundary Condition.
\end{proof}
\begin{definition}[Container Thermal Mass Model]
  \label{def:physics:phyx-mini-0406:phyxminiproblems-problemphyxmini0406-containerthermalmassmodel}
  \lean{PhyXMiniProblems.ProblemPhyXMini0406.ContainerThermalMassModel}
  Thermal-mass idealization for the beaker and gas container
\end{definition}
\begin{proof}
  This is the declared model definition of Container Thermal Mass Model.
\end{proof}
\begin{definition}[Figure Label]
  \label{def:physics:phyx-mini-0406:phyxminiproblems-problemphyxmini0406-figurelabel}
  \lean{PhyXMiniProblems.ProblemPhyXMini0406.FigureLabel}
  Literal labels visible in the supplied raster
\end{definition}
\begin{proof}
  This is the declared model definition of Figure Label.
\end{proof}
\begin{definition}[Figure Region]
  \label{def:physics:phyx-mini-0406:phyxminiproblems-problemphyxmini0406-figureregion}
  \lean{PhyXMiniProblems.ProblemPhyXMini0406.FigureRegion}
  The two material regions vertically separated by the thin metal line
\end{definition}
\begin{proof}
  This is the declared model definition of Figure Region.
\end{proof}
\begin{definition}[Figure Region Content]
  \label{def:physics:phyx-mini-0406:phyxminiproblems-problemphyxmini0406-figureregioncontent}
  \lean{PhyXMiniProblems.ProblemPhyXMini0406.FigureRegionContent}
  Material shown in a region of the figure
\end{definition}
\begin{proof}
  This is the declared model definition of Figure Region Content.
\end{proof}
\begin{definition}[Supplied Thermal Contact Figure]
  \label{def:physics:phyx-mini-0406:phyxminiproblems-problemphyxmini0406-suppliedthermalcontactfigure}
  \lean{PhyXMiniProblems.ProblemPhyXMini0406.SuppliedThermalContactFigure}
  \uses{def:physics:phyx-mini-0406:phyxminiproblems-problemphyxmini0406-figurelabel, def:physics:phyx-mini-0406:phyxminiproblems-problemphyxmini0406-figureregion, def:physics:phyx-mini-0406:phyxminiproblems-problemphyxmini0406-figureregioncontent}
  Qualitative information retained from the supplied figure
\end{definition}
\begin{proof}
  This is the declared model definition of Supplied Thermal Contact Figure.
\end{proof}
\begin{definition}[Gas State]
  \label{def:physics:phyx-mini-0406:phyxminiproblems-problemphyxmini0406-gasstate}
  \lean{PhyXMiniProblems.ProblemPhyXMini0406.GasState}
  \uses{def:physics:phyx-mini-0406:phyxminiproblems-problemphyxmini0406-volumequantity}
  A thermodynamic state of the monatomic gas
\end{definition}
\begin{proof}
  This is the declared model definition of Gas State.
\end{proof}
\begin{definition}[Water State]
  \label{def:physics:phyx-mini-0406:phyxminiproblems-problemphyxmini0406-waterstate}
  \lean{PhyXMiniProblems.ProblemPhyXMini0406.WaterState}
  A thermal state of the water
\end{definition}
\begin{proof}
  This is the declared model definition of Water State.
\end{proof}
\begin{definition}[Insulated Gas Water Setup]
  \label{def:physics:phyx-mini-0406:phyxminiproblems-problemphyxmini0406-insulatedgaswatersetup}
  \lean{PhyXMiniProblems.ProblemPhyXMini0406.InsulatedGasWaterSetup}
  \uses{def:physics:phyx-mini-0406:phyxminiproblems-problemphyxmini0406-massquantity, def:physics:phyx-mini-0406:phyxminiproblems-problemphyxmini0406-amountofsubstancescale, def:physics:phyx-mini-0406:phyxminiproblems-problemphyxmini0406-specificheatcapacityscale, def:physics:phyx-mini-0406:phyxminiproblems-problemphyxmini0406-thermalstage, def:physics:phyx-mini-0406:phyxminiproblems-problemphyxmini0406-gasconstitution, def:physics:phyx-mini-0406:phyxminiproblems-problemphyxmini0406-interfacematerial, def:physics:phyx-mini-0406:phyxminiproblems-problemphyxmini0406-thermalcontactquality, def:physics:phyx-mini-0406:phyxminiproblems-problemphyxmini0406-exteriorboundarycondition, def:physics:phyx-mini-0406:phyxminiproblems-problemphyxmini0406-containerthermalmassmodel, def:physics:phyx-mini-0406:phyxminiproblems-problemphyxmini0406-suppliedthermalcontactfigure, def:physics:phyx-mini-0406:phyxminiproblems-problemphyxmini0406-gasstate, def:physics:phyx-mini-0406:phyxminiproblems-problemphyxmini0406-waterstate}
  All independent physical quantities of the experiment. The final gas pressure is an unknown observable inside the final GasState; it is not defined from an answer choice or requested value
\end{definition}
\begin{proof}
  This is the declared model definition of Insulated Gas Water Setup.
\end{proof}
\begin{definition}[Matches Insulated Gas Water Scenario]
  \label{def:physics:phyx-mini-0406:phyxminiproblems-problemphyxmini0406-matchesinsulatedgaswaterscenario}
  \lean{PhyXMiniProblems.ProblemPhyXMini0406.MatchesInsulatedGasWaterScenario}
  \uses{def:physics:phyx-mini-0406:phyxminiproblems-problemphyxmini0406-amountofsubstancescale, def:physics:phyx-mini-0406:phyxminiproblems-problemphyxmini0406-specificheatcapacityscale, def:physics:phyx-mini-0406:phyxminiproblems-problemphyxmini0406-insulatedgaswatersetup}
  Qualitative apparatus and material assumptions stated in the prose
\end{definition}
\begin{proof}
  This is the declared model definition of Matches Insulated Gas Water Scenario.
\end{proof}
\begin{definition}[Matches Supplied Thermal Contact Figure]
  \label{def:physics:phyx-mini-0406:phyxminiproblems-problemphyxmini0406-matchessuppliedthermalcontactfigure}
  \lean{PhyXMiniProblems.ProblemPhyXMini0406.MatchesSuppliedThermalContactFigure}
  \uses{def:physics:phyx-mini-0406:phyxminiproblems-problemphyxmini0406-suppliedthermalcontactfigure}
  Qualitative labels and geometry read directly from the primary image
\end{definition}
\begin{proof}
  This is the declared model definition of Matches Supplied Thermal Contact Figure.
\end{proof}
\begin{definition}[Matches Problem Readouts]
  \label{def:physics:phyx-mini-0406:phyxminiproblems-problemphyxmini0406-matchesproblemreadouts}
  \lean{PhyXMiniProblems.ProblemPhyXMini0406.MatchesProblemReadouts}
  \uses{def:physics:phyx-mini-0406:phyxminiproblems-problemphyxmini0406-amountofsubstancescale, def:physics:phyx-mini-0406:phyxminiproblems-problemphyxmini0406-specificheatcapacityscale, def:physics:phyx-mini-0406:phyxminiproblems-problemphyxmini0406-massingrams, def:physics:phyx-mini-0406:phyxminiproblems-problemphyxmini0406-volumeincubiccentimeters, def:physics:phyx-mini-0406:phyxminiproblems-problemphyxmini0406-temperatureindegreescelsius, def:physics:phyx-mini-0406:phyxminiproblems-problemphyxmini0406-pressureinstandardatmospheres, def:physics:phyx-mini-0406:phyxminiproblems-problemphyxmini0406-insulatedgaswatersetup}
  Numerical data stated in the problem; no final pressure occurs here
\end{definition}
\begin{proof}
  This is the declared model definition of Matches Problem Readouts.
\end{proof}
\begin{definition}[Uses Reference Water Thermal Data]
  \label{def:physics:phyx-mini-0406:phyxminiproblems-problemphyxmini0406-usesreferencewaterthermaldata}
  \lean{PhyXMiniProblems.ProblemPhyXMini0406.UsesReferenceWaterThermalData}
  \uses{def:physics:phyx-mini-0406:phyxminiproblems-problemphyxmini0406-amountofsubstancescale, def:physics:phyx-mini-0406:phyxminiproblems-problemphyxmini0406-specificheatcapacityscale, def:physics:phyx-mini-0406:phyxminiproblems-problemphyxmini0406-insulatedgaswatersetup}
  Supplemental textbook thermal data needed for the numerical calculation. The problem expects this material property rather than printing it
\end{definition}
\begin{proof}
  This is the declared model definition of Uses Reference Water Thermal Data.
\end{proof}
\begin{definition}[Has Physical Thermodynamic Parameters]
  \label{def:physics:phyx-mini-0406:phyxminiproblems-problemphyxmini0406-hasphysicalthermodynamicparameters}
  \lean{PhyXMiniProblems.ProblemPhyXMini0406.HasPhysicalThermodynamicParameters}
  \uses{def:physics:phyx-mini-0406:phyxminiproblems-problemphyxmini0406-amountofsubstancescale, def:physics:phyx-mini-0406:phyxminiproblems-problemphyxmini0406-specificheatcapacityscale, def:physics:phyx-mini-0406:phyxminiproblems-problemphyxmini0406-pressureinpascals, def:physics:phyx-mini-0406:phyxminiproblems-problemphyxmini0406-massinkilograms, def:physics:phyx-mini-0406:phyxminiproblems-problemphyxmini0406-volumeincubicmeters, def:physics:phyx-mini-0406:phyxminiproblems-problemphyxmini0406-temperatureinkelvin, def:physics:phyx-mini-0406:phyxminiproblems-problemphyxmini0406-insulatedgaswatersetup}
  Positivity and nondegeneracy conditions for physical observables
\end{definition}
\begin{proof}
  This is the declared model definition of Has Physical Thermodynamic Parameters.
\end{proof}
\begin{definition}[Satisfies Insulated Gas Water Laws]
  \label{def:physics:phyx-mini-0406:phyxminiproblems-problemphyxmini0406-satisfiesinsulatedgaswaterlaws}
  \lean{PhyXMiniProblems.ProblemPhyXMini0406.SatisfiesInsulatedGasWaterLaws}
  \uses{def:physics:phyx-mini-0406:phyxminiproblems-problemphyxmini0406-amountofsubstancescale, def:physics:phyx-mini-0406:phyxminiproblems-problemphyxmini0406-specificheatcapacityscale, def:physics:phyx-mini-0406:phyxminiproblems-problemphyxmini0406-pressureinpascals, def:physics:phyx-mini-0406:phyxminiproblems-problemphyxmini0406-energyinjoules, def:physics:phyx-mini-0406:phyxminiproblems-problemphyxmini0406-massinkilograms, def:physics:phyx-mini-0406:phyxminiproblems-problemphyxmini0406-volumeincubicmeters, def:physics:phyx-mini-0406:phyxminiproblems-problemphyxmini0406-temperatureinkelvin, def:physics:phyx-mini-0406:phyxminiproblems-problemphyxmini0406-universalmolargasconstant-j-per-mol-k, def:physics:phyx-mini-0406:phyxminiproblems-problemphyxmini0406-insulatedgaswatersetup}
  Governing thermodynamic relations in coherent SI readouts: rigid volume, PV = nRT, monatomic-gas internal energy, water calorimetry, conservation of the isolated gas--water energy, and common long-time temperature. None contains a numerical final pressure or answer choice
\end{definition}
\begin{proof}
  This is the declared model definition of Satisfies Insulated Gas Water Laws.
\end{proof}
\begin{definition}[Answer Choice]
  \label{def:physics:phyx-mini-0406:phyxminiproblems-problemphyxmini0406-answerchoice}
  \lean{PhyXMiniProblems.ProblemPhyXMini0406.AnswerChoice}
  Labels of the four gas-pressure choices printed in the problem
\end{definition}
\begin{proof}
  This is the declared model definition of Answer Choice.
\end{proof}
\begin{definition}[displayed Pressure In Atmospheres]
  \label{def:physics:phyx-mini-0406:phyxminiproblems-problemphyxmini0406-displayedpressureinatmospheres}
  \lean{PhyXMiniProblems.ProblemPhyXMini0406.displayedPressureInAtmospheres}
  \uses{def:physics:phyx-mini-0406:phyxminiproblems-problemphyxmini0406-answerchoice}
  Pressure in standard atmospheres printed beside each answer label
\end{definition}
\begin{proof}
  This is the declared model definition of displayed Pressure In Atmospheres.
\end{proof}
\begin{definition}[recorded Dataset Answer]
  \label{def:physics:phyx-mini-0406:phyxminiproblems-problemphyxmini0406-recordeddatasetanswer}
  \lean{PhyXMiniProblems.ProblemPhyXMini0406.recordedDatasetAnswer}
  \uses{def:physics:phyx-mini-0406:phyxminiproblems-problemphyxmini0406-answerchoice}
  The dataset's recorded answer, retained only as metadata
\end{definition}
\begin{proof}
  This is the declared model definition of recorded Dataset Answer.
\end{proof}
\begin{definition}[Rounds To Displayed Pressure]
  \label{def:physics:phyx-mini-0406:phyxminiproblems-problemphyxmini0406-roundstodisplayedpressure}
  \lean{PhyXMiniProblems.ProblemPhyXMini0406.RoundsToDisplayedPressure}
  \uses{def:physics:phyx-mini-0406:phyxminiproblems-problemphyxmini0406-amountofsubstancescale, def:physics:phyx-mini-0406:phyxminiproblems-problemphyxmini0406-specificheatcapacityscale, def:physics:phyx-mini-0406:phyxminiproblems-problemphyxmini0406-pressureinstandardatmospheres, def:physics:phyx-mini-0406:phyxminiproblems-problemphyxmini0406-insulatedgaswatersetup, def:physics:phyx-mini-0406:phyxminiproblems-problemphyxmini0406-answerchoice, def:physics:phyx-mini-0406:phyxminiproblems-problemphyxmini0406-displayedpressureinatmospheres}
  A displayed pressure agrees after rounding to one decimal place
\end{definition}
\begin{proof}
  This is the declared model definition of Rounds To Displayed Pressure.
\end{proof}
\begin{definition}[Is Closest Displayed Pressure]
  \label{def:physics:phyx-mini-0406:phyxminiproblems-problemphyxmini0406-isclosestdisplayedpressure}
  \lean{PhyXMiniProblems.ProblemPhyXMini0406.IsClosestDisplayedPressure}
  \uses{def:physics:phyx-mini-0406:phyxminiproblems-problemphyxmini0406-amountofsubstancescale, def:physics:phyx-mini-0406:phyxminiproblems-problemphyxmini0406-specificheatcapacityscale, def:physics:phyx-mini-0406:phyxminiproblems-problemphyxmini0406-pressureinstandardatmospheres, def:physics:phyx-mini-0406:phyxminiproblems-problemphyxmini0406-insulatedgaswatersetup, def:physics:phyx-mini-0406:phyxminiproblems-problemphyxmini0406-answerchoice, def:physics:phyx-mini-0406:phyxminiproblems-problemphyxmini0406-displayedpressureinatmospheres}
  A choice is at least as close as every displayed pressure
\end{definition}
\begin{proof}
  This is the declared model definition of Is Closest Displayed Pressure.
\end{proof}
\begin{definition}[Is Unique Closest Displayed Pressure]
  \label{def:physics:phyx-mini-0406:phyxminiproblems-problemphyxmini0406-isuniqueclosestdisplayedpressure}
  \lean{PhyXMiniProblems.ProblemPhyXMini0406.IsUniqueClosestDisplayedPressure}
  \uses{def:physics:phyx-mini-0406:phyxminiproblems-problemphyxmini0406-amountofsubstancescale, def:physics:phyx-mini-0406:phyxminiproblems-problemphyxmini0406-specificheatcapacityscale, def:physics:phyx-mini-0406:phyxminiproblems-problemphyxmini0406-insulatedgaswatersetup, def:physics:phyx-mini-0406:phyxminiproblems-problemphyxmini0406-answerchoice, def:physics:phyx-mini-0406:phyxminiproblems-problemphyxmini0406-isclosestdisplayedpressure}
  The selected answer is the unique closest displayed pressure
\end{definition}
\begin{proof}
  This is the declared model definition of Is Unique Closest Displayed Pressure.
\end{proof}
\begin{lemma}[mass In Kilograms eq mass In Grams div thousand]
  \label{lem:physics:phyx-mini-0406:phyxminiproblems-problemphyxmini0406-massinkilograms-eq-massingrams-div-thousand}
  \lean{PhyXMiniProblems.ProblemPhyXMini0406.massInKilograms_eq_massInGrams_div_thousand}
  \uses{def:physics:phyx-mini-0406:phyxminiproblems-problemphyxmini0406-massquantity, def:physics:phyx-mini-0406:phyxminiproblems-problemphyxmini0406-massingrams, def:physics:phyx-mini-0406:phyxminiproblems-problemphyxmini0406-massinkilograms}
  Unit-conversion consequences of the dimensional quantity types
\end{lemma}
\begin{proof}
  The conclusion follows from the governing relations and figure readouts named in the statement of mass In Kilograms eq mass In Grams div thousand.
\end{proof}
\begin{lemma}[volume In Cubic Meters eq volume In Cubic Centimeters div million]
  \label{lem:physics:phyx-mini-0406:phyxminiproblems-problemphyxmini0406-volumeincubicmeters-eq-volumeincubiccentimeters-div-million}
  \lean{PhyXMiniProblems.ProblemPhyXMini0406.volumeInCubicMeters_eq_volumeInCubicCentimeters_div_million}
  \uses{def:physics:phyx-mini-0406:phyxminiproblems-problemphyxmini0406-volumequantity, def:physics:phyx-mini-0406:phyxminiproblems-problemphyxmini0406-volumeincubiccentimeters, def:physics:phyx-mini-0406:phyxminiproblems-problemphyxmini0406-volumeincubicmeters}
  This lemma records the volume In Cubic Meters eq volume In Cubic Centimeters div million component of the problem-specific physical model
\end{lemma}
\begin{proof}
  The conclusion follows from the governing relations and figure readouts named in the statement of volume In Cubic Meters eq volume In Cubic Centimeters div million.
\end{proof}
\begin{lemma}[pressure In Pascals eq atmospheres mul standard Atmosphere]
  \label{lem:physics:phyx-mini-0406:phyxminiproblems-problemphyxmini0406-pressureinpascals-eq-atmospheres-mul-standardatmosphere}
  \lean{PhyXMiniProblems.ProblemPhyXMini0406.pressureInPascals_eq_atmospheres_mul_standardAtmosphere}
  \uses{def:physics:phyx-mini-0406:phyxminiproblems-problemphyxmini0406-pressureinpascals, def:physics:phyx-mini-0406:phyxminiproblems-problemphyxmini0406-pressureinstandardatmospheres}
  This lemma records the pressure In Pascals eq atmospheres mul standard Atmosphere component of the problem-specific physical model
\end{lemma}
\begin{proof}
  The conclusion follows from the governing relations and figure readouts named in the statement of pressure In Pascals eq atmospheres mul standard Atmosphere.
\end{proof}
\begin{lemma}[final Equilibrium Temperature In Kelvin eq]
  \label{lem:physics:phyx-mini-0406:phyxminiproblems-problemphyxmini0406-finalequilibriumtemperatureinkelvin-eq}
  \lean{PhyXMiniProblems.ProblemPhyXMini0406.finalEquilibriumTemperatureInKelvin_eq}
  \uses{def:physics:phyx-mini-0406:phyxminiproblems-problemphyxmini0406-amountofsubstancescale, def:physics:phyx-mini-0406:phyxminiproblems-problemphyxmini0406-specificheatcapacityscale, def:physics:phyx-mini-0406:phyxminiproblems-problemphyxmini0406-temperatureinkelvin, def:physics:phyx-mini-0406:phyxminiproblems-problemphyxmini0406-insulatedgaswatersetup, def:physics:phyx-mini-0406:phyxminiproblems-problemphyxmini0406-matchesproblemreadouts, def:physics:phyx-mini-0406:phyxminiproblems-problemphyxmini0406-usesreferencewaterthermaldata, def:physics:phyx-mini-0406:phyxminiproblems-problemphyxmini0406-satisfiesinsulatedgaswaterlaws}
  The energy balance determines a common equilibrium temperature of about 345.20 K; the exact rational SI readout is retained here
\end{lemma}
\begin{proof}
  The conclusion follows from the governing relations and figure readouts named in the statement of final Equilibrium Temperature In Kelvin eq.
\end{proof}
\begin{lemma}[final Gas Pressure In Standard Atmospheres eq]
  \label{lem:physics:phyx-mini-0406:phyxminiproblems-problemphyxmini0406-finalgaspressureinstandardatmospheres-eq}
  \lean{PhyXMiniProblems.ProblemPhyXMini0406.finalGasPressureInStandardAtmospheres_eq}
  \uses{def:physics:phyx-mini-0406:phyxminiproblems-problemphyxmini0406-amountofsubstancescale, def:physics:phyx-mini-0406:phyxminiproblems-problemphyxmini0406-specificheatcapacityscale, def:physics:phyx-mini-0406:phyxminiproblems-problemphyxmini0406-pressureinstandardatmospheres, def:physics:phyx-mini-0406:phyxminiproblems-problemphyxmini0406-insulatedgaswatersetup, def:physics:phyx-mini-0406:phyxminiproblems-problemphyxmini0406-matchesproblemreadouts, def:physics:phyx-mini-0406:phyxminiproblems-problemphyxmini0406-usesreferencewaterthermaldata, def:physics:phyx-mini-0406:phyxminiproblems-problemphyxmini0406-satisfiesinsulatedgaswaterlaws}
  The final ideal-gas equation then gives approximately 2.83245 atm
\end{lemma}
\begin{proof}
  The conclusion follows from the governing relations and figure readouts named in the statement of final Gas Pressure In Standard Atmospheres eq.
\end{proof}
% --- Archon declaration topology end ---
% --- Archon physics formalization source end ---
